% ===== §3 The Ontology of the Field =====
\section{The Ontology of the Field}
\label{p22:sec:field-ontology}

\noindent
The reduction chain that organises this paper begins from a claim carried over from the immediately preceding work, that relational production requires a field. The companion paper established relational production as the process by which the relational subject is generated, and it established that this generation does not occur in a void but within a condition, an openness it named under the figure of spaciousness. This section takes up that condition and gives it a general form, because the chain's first link, relational production requires a field, needs the field to be specified before the links that follow can carry weight. What is a field, such that relational production requires one, such that its construction will turn out to be an aesthetic rather than an engineering matter? The answer given here is that a field is a conditional environment of coupling, and the section's work is to make that definition precise, to distinguish it from the physical space that is only its carrier, and to relate it to the spaciousness from which it descends without collapsing into it.

\subsection{The field: a conditional environment of coupling}
\label{p22:sub:field-definition}

A field, in the sense this paper requires, is the conditional structure within which two predictive-valuing systems can couple in the generative direction. It is not a place but a condition, or more exactly it is the set of conditions under which the coupling that produces relational value can occur and accumulate: the rhythm of a shared life, its atmosphere, the distribution of attention within it, the arrangement of things insofar as that arrangement bears on whether the coupling gains or spins. These conditions are what a field is, and they are conditions of possibility for accumulation, not the accumulation itself and not the subjects that accumulate. A physical space, a kitchen or a room, is a field only derivatively, as the carrier of such conditions, the place in which the rhythm is kept and the attention distributed and the things arranged; the space is where the field is held, but the field is the conditional structure the space carries, and the same space carries different fields as its conditions change, while the same field can be carried, imperfectly, across different spaces.

\claim{A field is a conditional environment of coupling: the structure of conditions under which two predictive-valuing systems can couple in the generative direction and accumulate holonomy. It is not a physical space but a condition that physical space may carry; rhythm, atmosphere, the distribution of attention, and the arrangement of things are conditions of the field insofar as they bear on whether coupling gains or spins. The field is therefore neither the subjects nor their accumulated value but the conditional structure that makes their generative coupling possible, and to build a field is to shape those conditions.}

The reason the definition must be pitched at the level of conditions, and not of physical space, is that everything in the chain depends on it, and the dependence can be stated now even though its payoff comes later. If the field were simply physical space, its construction would indeed be an engineering matter, since physical space can be designed, optimised, and standardised, and the chain's second link, that field-building is aesthetic rather than engineering, would fail at once. It is because the field is a structure of conditions, whose appropriateness is state-dependent and cannot be standardised, that its construction escapes engineering, and the next section's whole argument rests on the field being conditional in this way. The ontology is thus not a preliminary nicety but the load-bearing specification on which the chain's most contested link stands or falls, and it is set out here so that the link can be argued there without re-litigating what a field is.

\begin{casebox}[The same room, and the field that is not the room]
A kitchen is a kitchen, the same cabinets and the same light, and yet the field it carries is not constant. On one evening the kitchen holds a generative field, the two moving in a rhythm that gains, the room thick with a coupling that accumulates; on another, after a quarrel, the identical kitchen holds a flattened field, the same cabinets and light carrying now a silence that spins. The room has not changed; the field has, which shows the field to be not the room but a condition the room carries, able to change while the room stays fixed. And the converse holds: a couple who move to a new house carry the field with them, imperfectly, the rhythm and atmosphere and distribution of attention re-established in a different space, so that the same field, or nearly the same, is borne by a new room. The field travels though the room cannot, and the room changes its field though it cannot change itself, and these two ordinary facts together show what the definition asserts, that the field is the conditional structure and the physical space only its carrier, the two able to vary independently because they are not the same thing.
\end{casebox}

\subsection{The field and the subjects it couples}
\label{p22:sub:field-subjects}

A field couples subjects, and the relation between the field and the subjects it couples must be stated carefully to avoid two errors. The first error would make the field a third thing alongside the two subjects, a container in which they sit, as though the conditions of coupling existed independently of the coupling they condition. They do not. The conditions of a field are sustained by the very coupling they make possible, in the way that the rhythm of a shared life is kept by the partners whose sharing it conditions, so that the field is not a container prior to its occupants but a structure they hold up by inhabiting it, a condition that is maintained in being by what it conditions. This circularity is not vicious; it is the ordinary structure of a sustained relation, in which the conditions of coupling and the coupling are mutually holding, each keeping the other in existence, and it is why a field can decay when the coupling lapses and can be rebuilt when the coupling resumes.

The second error would make the field nothing over the subjects, a mere name for their interaction, with no structure of its own to be built or damaged. This too is false, because the conditions of coupling have a structure that outruns any single interaction and shapes interactions yet to come: a field whose rhythm has been established, whose atmosphere has been tended, whose attention is well distributed, makes generative coupling easier for the next interaction, and a field whose conditions have been damaged makes it harder, regardless of the good will brought to the particular moment. The field is thus real enough to be built and damaged, to make coupling easier or harder, to be the standing result of past coupling that conditions future coupling, while not being a thing apart from the subjects whose coupling sustains it. It is, in the framework's terms, the accumulated condition of a relation, the settled structure within which the next cycle of perception, reception, sharing, and reproduction will be easier or harder to turn, and this is exactly why its construction is the proper object of an aesthetic education: the field is what the cultivation builds.

\subsection{The field and spaciousness}
\label{p22:sub:field-spaciousness}

The field descends from the spaciousness of the companion work but is more general than it, and the relation should be made exact so that this paper extends rather than repeats. Spaciousness was the openness, the room left for what is not yet determined, the condition under which contingency could be met and the relational subject generated in the field of travel; it was a particular condition of the field, the condition of openness, given prominence because travel is the domain in which openness is most plainly the generative condition. The field, as this paper defines it, is the more general structure of which spaciousness is one possible state. A field can be open, spacious, full of room for contingency, and then spaciousness is its condition; but a field can also be settled, dense, repetitive, with little room for the new, and such a field is not therefore no field, but a field in a different state, whose generativity must be secured by means other than openness. This generality is what lets the present paper move where the companion could not, into the dense and repetitive and unspacious conditions of the most ordinary shared life, the conditions under which there is no room to travel and the field must accumulate by repetition rather than by openness. The field is the genus; spaciousness is one of its states; and the cultivation this paper describes is the building of generative fields across all their states, including the dense and repetitive ones in which spaciousness is not available and the good cycle must be turned in the narrow room that the everyday affords.

\subsection{The depth of a field}
\label{p22:sub:field-depth}

A field has not only a direction, whether it gains or flattens, but a depth, and introducing the notion completes the ontology in a way the later sections will draw on. The depth of a field is the magnitude of the accumulation it can sustain, the difference between a shallow field that gains a little and a deep one that gains greatly, between a coupling that accumulates a thin holonomy and one that accumulates a thick one. Two fields may both be generative, both gaining rather than flattening, and yet differ in depth, the one a pleasant and slight accumulation, the other a profound one, and the difference is not in the direction, which is the same, but in the magnitude, the depth of the gain. This is why a relation can be good, genuinely accumulating, and yet shallow, and why the deepening of a field is a distinct achievement from the establishing of its generative direction; one first builds a field that gains, and then, over time, deepens the gain, the cultivation of maintenance being in large part the deepening of a field already generative, the thickening of an accumulation already underway.

The depth of a field is built by time and by the four-step cycle turned many times, and this connects the ontology to the account of reproduction. A field deepens as its cycle is turned, each turn that gains adding to the depth, so that the deep field is the one whose cycle has been turned many times well, the accumulation thickened by long good repetition. This is why depth cannot be rushed and why the deep field is characteristically the long one, the relation that has accumulated over years of good cycles, its depth the sediment of a thousand small gains. It is also why the loss of a deep field is felt as a greater loss than the loss of a shallow one, since what is lost is not only a present accumulation but the depth built by all the past cycles that thickened it, the long accumulation that cannot be quickly rebuilt because it was the work of long repetition. The depth of a field is thus the dimension along which the temporal achievement of maintenance registers, the thickness of accumulation that long good repetition builds, and the ontology must include it because the cultivation aims not only at the generative direction of a field but at its depth, not only at building a field that gains but at deepening the gain over the time that a shared life affords.
